\section*{第 \ref{ch:g12:special-relativity} 章——狭义相对论：时间延缓}

\begin{solution}{exo:g12:special-relativity:1}
$\gamma = 1/\sqrt{1 - (v/c)^2}$：$1.01$（$0.1c$）；$1.25$（$0.6c$）；
$1.67$（$0.8c$）；$10.0$（$0.995c$）。
\end{solution}

\begin{solution}{exo:g12:special-relativity:2}
$\gamma = 1.25$，故地面测得 $1.25 \times 10.0 =
\qty{12.5}{s}$。飞船上的 \qty{10.0}{s} 是固有时：
节拍器（同一只钟）在两次拍击时都在场。
\end{solution}

\begin{solution}{exo:g12:special-relativity:3}
$v = c\sqrt{1 - 1/\gamma^2}$：$\gamma = 2$ 时 $v = 0.866c$；
$\gamma = 10$ 时 $v = 0.995c$。
\end{solution}

\begin{solution}{exo:g12:special-relativity:4}
$v = \qty{36.1}{m/s}$，$v/c = \num{1.2e-7}$；
$\gamma - 1 \approx (\num{1.2e-7})^2/2 \approx \num{7.2e-15}$。损失一秒需要
$1/\num{7.2e-15} \approx \qty{1.4e14}{s}$ ——约
四百万年的驾驶时间。
\end{solution}

\begin{solution}{exo:g12:special-relativity:5}
乘客迎着前端闪光运动，因此闪光先到达他处；
由第二假设，两道闪光以相同的 $c$ 传播，经过
在\emph{他的参考系中}相等的半车厢距离，因此
先到达即先发生。同时性与参考系有关：每个
结论在其各自的参考系中都是正确的。
\end{solution}

\begin{solution}{exo:g12:special-relativity:6}
$\gamma = 1/\sqrt{1 - 0.99^2} = 7.09$；实验室中的寿命为
$7.09 \times \num{2.6e-8} = \qty{1.8e-7}{s}$。经典射程
$0.99 \times \num{3.00e8} \times \num{2.6e-8} \approx \qty{7.7}{m}$；
实际射程 $7.09 \times 7.7 \approx \qty{55}{m}$。
\end{solution}

\begin{solution}{exo:g12:special-relativity:7}
当 $\gamma = 1.5$ 时 $v = c\sqrt{1 - 1/2.25} \approx 0.75c$；
$\gamma = 3$ 时 $\approx 0.94c$。渐近情形：$\gamma$（以及
继续加速所需的能量代价）随 $v \to c$ 趋于无穷——
无论付出多大努力，都无法使有质量的物体达到光速。
\end{solution}

\begin{solution}{exo:g12:special-relativity:8}
$\Delta t_0 = 2L/c = 3.0/\num{3.00e8} = \qty{1.0e-8}{s}$；在 $0.8c$ 时，
$\gamma = 5/3$，故 $\Delta t = \qty{1.7e-8}{s}$。
\end{solution}

\begin{solution}{exo:g12:special-relativity:9}
固有时属于在两个事件处都在场的钟：(a) 佩戴者；
(b) μ 子；(c) 飞行员——只有飞船同时在
出发与到达两个事件处都在场。
\end{solution}

\begin{solution}{exo:g12:special-relativity:10}
$\gamma - 1 \approx 7700^2/(2 \times \num{9.0e16}) \approx
\num{3.3e-10}$；经过 \qty{1.58e7}{s} 后，宇航员落后
$\num{3.3e-10} \times \num{1.58e7} \approx \qty{5.2e-3}{s}$ ——约
年轻五毫秒。
\end{solution}

\begin{solution}{exo:g12:special-relativity:11}
$E = \num{1.0e-3} \times (\num{3.00e8})^2 = \qty{9.0e13}{J}$。
一座 \qty{1.0}{GW} 的电站每天输出 $\num{1.0e9} \times 86400 \approx
\qty{8.6e13}{J}$ ——一克物质相当于该电站一天的输出，而
$\num{9.0e13}/\num{8.5e5} \approx \num{e8}$：相当于一亿辆
飞驰的汽车。
\end{solution}

\begin{solution}{exo:g12:special-relativity:12}
两边平方并取倒数：$\gamma^2(1 - v^2/c^2) = 1$，故
$v^2/c^2 = 1 - 1/\gamma^2$，因此 $v = c\sqrt{1 - 1/\gamma^2}$。
当 $\gamma = 30$ 时，$v = c\sqrt{1 - 1/900} = 0.99944c$。
\end{solution}

\begin{solution}{exo:g12:special-relativity:13}
收缩后的厚度 $L = \qty{15}{km}/30 = \qty{500}{m}$；穿越时间
$500/(0.99944 \times \num{3.00e8}) \approx \qty{1.7e-6}{s} <
\qty{2.2}{\micro s}$。两个参考系都同意 μ 子能到达地面——
一方归因于寿命被拉长，另一方归因于大气层被压缩。
\end{solution}

\begin{solution}{exo:g12:special-relativity:14}
实验室中：$3200/\num{3.00e8} \approx \qty{1.1e-5}{s}$（速率
$\approx c$）。电子参考系中：管长收缩为
$3200/\num{1.0e5} = \qty{3.2}{cm}$，穿越时间为
$\num{1.1e-5}/\num{1.0e5} \approx \qty{1.1e-10}{s}$。
\end{solution}

\begin{solution}{exo:g12:special-relativity:15}
地球上：$4.2/0.9 \approx \qty{4.7}{years}$。由于
$\gamma = 1/\sqrt{1 - 0.81} = 2.29$，探测器的钟记录
$4.7/2.29 \approx \qty{2.0}{years}$。
\end{solution}

\begin{solution}{pb:g12:special-relativity:1}
\textbf{1.} $c\,\Delta t_0 = \num{3.00e8} \times \num{2.2e-6} \approx
\qty{660}{m}$。

\textbf{2.} $\num{15000}/660 \approx 23$ 个寿命。

\textbf{3.} $\mathrm{e}^{-23} \approx \num{e-10}$：每一百亿个 μ 子中
应只有不到一个到达，然而通量是每 \unit{cm^2} 每
分钟一个——经典物理被天空彻底否定了。

\textbf{4.} 地面寿命 $30 \times \qty{2.2}{\micro s} =
\qty{66}{\micro s}$；射程 $\approx \num{3.00e8} \times \num{6.6e-5}
\approx \qty{20}{km} > \qty{15}{km}$：到达现象得到了解释。

\textbf{5.} $v = c\sqrt{1 - 1/900} = 0.99944c$.

\textbf{6.} $\Delta t_0 = 2L/c = 3.0/\num{3.00e8} = \qty{1.0e-8}{s}$.

\textbf{7.} 水平方向为 $v\,\Delta t/2$；竖直方向为 $L$；倾斜的
半程长度为 $c\,\Delta t/2$，因为光速不变原理确定了
光子在地面参考系中的速率同样为 $c$。

\textbf{8.} $(c\,\Delta t/2)^2 = L^2 + (v\,\Delta t/2)^2$ 可得
$\Delta t^2(c^2 - v^2) = 4L^2$，于是 $\Delta t =
(2L/c)/\sqrt{1 - v^2/c^2} = \gamma\,\Delta t_0$.

\textbf{9.} $\gamma = 1/\sqrt{1 - 0.36} = 1.25$;
$\Delta t = \qty{1.25e-8}{s}$.

\textbf{10.} 飞船上的：其唯一的钟在两次滴答时都在场，因此它
读出的是固有时；地面则需要两只同步的钟。

\textbf{11.} $\gamma = 1/\sqrt{1 - 0.64} = 1/0.6 = 5/3 \approx 1.67$.

\textbf{12.} $10.0/(5/3) = \qty{6.0}{years}$.

\textbf{13.} 去程 $\qty{5.0}{years}$，以 $0.8c$ 运动：$4.0$ 光年。

\textbf{14.} $10.0 - 6.0 = \qty{4.0}{years}$：宇航员比
她的孪生同胞年轻 4 岁。

\textbf{15.} 两者的处境并不对称：只有宇航员掉头
——她加速并改变惯性参考系，而地球上的
孪生同胞始终处于同一惯性参考系中，因此只有她的钟记录了较短的时间。

\textbf{16.} $v = 2\pi \times \num{2.66e7}/\num{43200} \approx
\qty{3.9e3}{m/s}$.

\textbf{17.} $\gamma - 1 \approx (\num{3.87e3})^2/(2 \times
\num{9.0e16}) \approx \num{8.3e-11}$.

\textbf{18.} $\num{8.3e-11} \times \num{86400} \approx \qty{7.2e-6}{s}
= \qty{7.2}{\micro s}$ 每天慢。

\textbf{19.} $c \times \num{7.2e-6} \approx \qty{2.2}{km}$，即每天的
位置误差。

\textbf{20.} 净漂移 $45 - 7.2 \approx \qty{38}{\micro s}$，
每天快；误差 $c \times \num{38e-6} \approx \qty{11}{km}$，每天——
这正是每台 GPS 钟预先支付的相对论代价：地面上预先调慢，
使其在轨道上走时准确。
\end{solution}
